\documentclass[a4paper,11pt]{article}
\usepackage[T2A]{fontenc}
\usepackage[utf8]{inputenc}
\usepackage[russian]{babel}
\usepackage{amsmath,amssymb,mathtools}
\usepackage{geometry}
\geometry{top=20mm,bottom=20mm,left=21mm,right=21mm}
\usepackage{microtype}
\usepackage{booktabs,tabularx,array}
\usepackage{enumitem}
\usepackage{xcolor}
\usepackage{tikz}
\usepackage{fancyhdr}
\usepackage{hyperref}
\usepackage{parskip}
\definecolor{accent}{HTML}{087C85}
\definecolor{deep}{HTML}{16585F}
\definecolor{error}{HTML}{A33C37}
\definecolor{body}{HTML}{26292A}
\hypersetup{colorlinks=true,linkcolor=deep,bookmarksopen=true}
\setlength{\parindent}{0pt}
\setlength{\parskip}{4pt}
\linespread{1.055}
\setlist[enumerate]{leftmargin=*,itemsep=3pt,topsep=4pt}
\pagestyle{fancy}
\fancyhf{}
\renewcommand{\headrulewidth}{0pt}
\fancyfoot[C]{\footnotesize\color{deep}Лёвин Артём Александрович эксклюзивно для Анастасии Павловой}
\fancyfoot[R]{\footnotesize\thepage}
\setlength{\footskip}{18pt}
\newcommand{\wrong}{\textcolor{error}{\textbf{Ошибка:}}}
\newcommand{\idea}{\textcolor{deep}{\textbf{Ключевая идея:}}}
\newcommand{\result}{\textcolor{accent}{\textbf{Итог:}}}
\newcommand{\tasktitle}[1]{\phantomsection\label{task:#1}{\Large\color{deep}\textbf{Задание №\,#1}}\par\vspace{4mm}}
\begin{document}
\color{body}
\thispagestyle{empty}
\begin{center}
\vspace*{27mm}
{\small\color{deep}\MakeUppercase{Математика}}\par
\vspace{20mm}
{\LARGE Ревью домашней работы}\par
\vspace{12mm}
{\Large\color{accent}Кредиты, вклады и годовые проценты}\par
\vspace{12mm}
{\normalsize Проверка решений и подробный разбор ошибок}\par
\vfill
{\large Анастасия Павлова}\par
\vspace{9mm}
{8 ноября 2026 г.}\par
\vspace{15mm}
{\small Лёвин Артём Александрович}\par
{\small эксклюзивно для Анастасии Павловой}\par
\vspace{15mm}
\begin{tikzpicture}[x=1cm,y=1cm]
\draw[accent,line width=0.7pt] (0,0) .. controls (1.8,0.6) and (2.7,-0.5) .. (4.3,0.0)
 .. controls (5.5,0.55) and (6.7,-0.5) .. (8.2,0);
\end{tikzpicture}
\end{center}

\clearpage
{\Large\color{deep}\textbf{Сводная таблица}}\par
\vspace{2mm}
В документ включены задания, в которых необходимо исправить рассуждение,
уточнить расчёт или дополнить доказательство. Правильно выполненные задания не приводятся.

\vspace{5mm}
\small
\renewcommand{\arraystretch}{1.38}
\begin{tabularx}{\linewidth}{@{}p{12mm}p{25mm}X p{26mm}p{28mm}@{}}
\toprule
№ & Тема & Суть недочёта & Ваш ответ & Правильный ответ\\
\midrule
\hyperref[task:2]{2} & Схемы кредита &
В пункте «в» требуется обоснование классификации. &
«Не относится»; пояснение не раскрыто. &
Не является ни одной из двух стандартных схем.\\
\addlinespace
\hyperref[task:7]{7} & Платежи по годам &
Проценты второго и третьего года рассчитаны в обратном порядке. &
354\,000; 318\,000; 336\,000 руб. &
354\,000; 336\,000; 318\,000 руб.\\
\addlinespace
\hyperref[task:8]{8} & Разность платежей &
Последний платёж рассчитан по неверной формуле; доказательство для соседних лет отсутствует. &
На 10\,000 руб.; без обоснования. &
Каждый следующий платёж меньше на 10\,000 руб.\\
\addlinespace
\hyperref[task:9]{9} & Целая ставка &
Не завершён отбор допустимой целой процентной ставки. &
Указаны 7\% и 7{,}7\% с вопросом. &
7\%; общая выплата 1\,210\,000 руб.\\
\bottomrule
\end{tabularx}
\normalsize

\vspace{10mm}
{\color{deep}\textbf{Общий ориентир.}} При равном погашении основного долга
проценты каждый год начисляются на \emph{остаток задолженности в начале года}.
В задачах с ограничениями на ставку нужно отдельно учитывать, какие значения
ставки разрешены условием.

\clearpage
\thispagestyle{empty}
\vspace*{\fill}
\begin{center}
{\LARGE\color{deep}\textbf{Разбор ошибок}}
\end{center}
\vspace*{\fill}

\clearpage
\tasktitle{2}
{\textbf{Условие.}} Определите схему погашения: 
а) все годовые платежи равны;
б) ежегодно погашается одинаковая часть основного долга;
в) остаток основного долга ежегодно уменьшается на 25\% от \emph{текущего остатка}.
Поясните, относится ли пункт «в» сам по себе к одной из двух стандартных схем.

\textbf{Ваше решение.} Для пункта «а» указана аннуитетная схема, для пункта «б» ---
дифференцированная. В пункте «в» отмечено, что это не одна из двух стандартных
схем, однако математическая причина не раскрыта.

\textbf{Ваш ответ.} Классификация по существу верная; обоснование пункта «в»
недостаточно.

\wrong\ Не хватает проверки определяющего признака схемы.

\textbf{Где именно возникает недочёт.} Последний правильный шаг ---
распознавание того, что в пункте «в» речь идёт о проценте от \emph{текущего}
остатка. Следующий необходимый шаг --- вычислить, как изменяется сумма
погашаемого основного долга; этого шага в записи нет.

\textbf{Почему этого недостаточно.} Для дифференцированной схемы погашение
\emph{основного долга} одинаково каждый год. Для аннуитетной схемы одинаков
\emph{полный платёж}, включающий проценты и погашение долга. В пункте «в»
задаётся иной закон уменьшения задолженности. Чтобы подтвердить вывод,
нужно проверить оба признака.

\idea\ Если каждый год остаётся 75\% предыдущего остатка, задолженность
образует геометрическую последовательность.

\textbf{Исправление и полное решение.} Пусть первоначальный долг равен
\(S>0\). После первого года остаток:
\[
B_1=S-0{,}25S=0{,}75S.
\]
За первый год погашено \(0{,}25S\) основного долга. После второго года:
\[
B_2=0{,}75B_1=0{,}75^2S,
\]
а погашено
\[
B_1-B_2=0{,}25\cdot0{,}75S=0{,}1875S.
\]
Таким образом, части основного долга \emph{разные}:
\(0{,}25S\ne0{,}1875S\).
Значит, схема не является дифференцированной.

Если ставка за год равна \(p\ge0\) (доля единицы), то первый полный платёж
составляет \((0{,}25+p)S\), а второй ---
\[
(0{,}25+p)\,0{,}75S.
\]
Второй платёж равен 75\% первого. Следовательно, он также не постоянен:
схема не является аннуитетной.

\textbf{Проверка.} При \(S=1000\) руб. за первый год возвращается
250 руб. долга, за второй --- 187{,}50 руб. долга. Равенства частей нет.

\result\ Пункт «а» --- аннуитетная схема; пункт «б» ---
дифференцированная; пункт «в» --- ни одна из них.

\textbf{Задача на закрепление.} Остаток долга в конце каждого года
уменьшается на 20\% от остатка в начале этого года. Объясните с помощью
двух первых лет, является ли такая схема аннуитетной или
дифференцированной.

\clearpage
\tasktitle{7}
{\textbf{Условие.}} Кредит 900\,000 руб. взят на три года под 6\% годовых.
Каждый год погашается одинаковая часть основного долга.
Найдите три годовых платежа, общую выплату и переплату.

\textbf{Ваше решение.} Верно найдены общая выплата 1\,008\,000 руб.
и переплата 108\,000 руб. Первый платёж вычислен как 354\,000 руб.
Затем записаны платежи 318\,000 руб. и 336\,000 руб.

\textbf{Ваш ответ.} По годам:
354\,000; 318\,000; 336\,000 руб.;
всего 1\,008\,000 руб.; переплата 108\,000 руб.

\wrong\ Перепутан остаток долга, на который начисляются проценты
во втором и третьем годах.

\textbf{Где именно возникает ошибка.} Последний правильный шаг ---
первый платёж:
\[
\frac{900\,000}{3}+900\,000\cdot0{,}06
=300\,000+54\,000=354\,000.
\]
Первый неправильный переход --- вычисление \emph{второго} платежа
по формуле
\[
\frac{900\,000}{3}+\frac{900\,000}{3}\cdot0{,}06
=318\,000.
\]
Здесь использован остаток 300\,000 руб. Он возникает только
к началу третьего года.

\textbf{Почему это неверно.} При дифференцированном погашении ежегодно
возвращается одна и та же часть тела кредита:
\[
900\,000:3=300\,000\ \text{руб.}
\]
После первого платежа остаётся \(900\,000-300\,000=600\,000\) руб.
Именно на эти 600\,000 руб. начисляются проценты за второй год.
Третий год начинается с остатком 300\,000 руб.

\idea\ Для каждого года используйте правило
\[
\text{Платёж}=
\text{часть основного долга}+
\text{ставка}\cdot\text{долг на начало года}.
\]

\newpage
\textbf{Исправление.} Второй платёж:
\[
300\,000+600\,000\cdot0{,}06
=300\,000+36\,000=\boxed{336\,000}.
\]
Третий платёж:
\[
300\,000+300\,000\cdot0{,}06
=300\,000+18\,000=\boxed{318\,000}.
\]

\textbf{Полное правильное решение.}
\begin{center}
\small
\begin{tabular}{@{}c r r r@{}}
\toprule
Год & Долг в начале, руб. & Проценты, руб. & Платёж, руб.\\
\midrule
1 & 900\,000 & 54\,000 & 354\,000\\
2 & 600\,000 & 36\,000 & 336\,000\\
3 & 300\,000 & 18\,000 & 318\,000\\
\bottomrule
\end{tabular}
\end{center}
Общая выплата:
\[
354\,000+336\,000+318\,000=\boxed{1\,008\,000\ \text{руб.}}
\]
Переплата:
\[
1\,008\,000-900\,000=\boxed{108\,000\ \text{руб.}}
\]

\textbf{Проверка.} Сумма погашенного основного долга:
\(3\cdot300\,000=900\,000\) руб.
Сумма процентов:
\(54\,000+36\,000+18\,000=108\,000\) руб.
Итоги совпадают. Платежи каждый год уменьшаются на
\(18\,000\) руб., что согласуется с убыванием остатка.

\result\ Платежи по годам: \textbf{354\,000, 336\,000, 318\,000 руб.}
Общая выплата и переплата были найдены верно.

\textbf{Задача на закрепление.} Кредит 600\,000 руб. взят на три года
под 8\% годовых. Основной долг погашают равными частями в конце каждого
года. Найдите все три годовых платежа, их сумму и переплату.

\clearpage
\tasktitle{8}
{\textbf{Условие.}} Кредит 800\,000 руб. взят на четыре года под 5\% годовых.
Основной долг погашается равными частями. На сколько рублей
каждый следующий ежегодный платёж меньше предыдущего?

\textbf{Ваше решение.} Верно вычислен первый платёж 240\,000 руб. Далее записано
\[
\frac{800\,000}{4}+\frac{200\,000}{4}\cdot0{,}05=210\,000,
\]
что неверно как арифметическое равенство. Число 210\,000 руб. для
последнего платежа всё же правильное. Указано «на 10\,000 рублей»
с замечанием, что доказательство пока не найдено.

\textbf{Ваш ответ.} На 10\,000 руб. Число верное.

\wrong\ Неверно записаны проценты в четвёртом платеже и отсутствует
обоснование разности \emph{соседних} платежей.

\textbf{Где именно возникает недочёт.} Последний правильный шаг --- вычисление первого платежа. Первый
неверный шаг --- выражение для четвёртого платежа:
\[
\frac{800\,000}{4}+\frac{200\,000}{4}\cdot0{,}05
=202\,500\ne210\,000.
\]
В числителе процентов ошибочно появляется дополнительное деление
на четыре. Далее также пропущено доказательство одинакового
уменьшения именно соседних платежей.

\textbf{Почему это важно.} Проценты в последнем году начисляются на остаток
200\,000 руб. и составляют 10\,000 руб. Поэтому корректная запись
четвёртого платежа:
\[
200\,000+200\,000\cdot0{,}05=210\,000.
\]
Разность между первым и последним платежами составляет \(240\,000-210\,000=30\,000\) руб.
Между четырьмя платежами есть \emph{три} перехода к следующему году.
Именно поэтому разность соседних платежей может быть 10\,000 руб.,
однако это ещё нужно доказать из условий кредита.

\idea\ Сумма ежегодно погашаемого долга постоянна, поэтому изменение
платежа определяется только уменьшением начисленных процентов.

\newpage
\textbf{Исправление и полное решение.}
Часть основного долга, возвращаемая ежегодно:
\[
d=\frac{800\,000}{4}=200\,000\ \text{руб.}
\]
Остатки в начале четырёх лет:
\[
800\,000;\quad600\,000;\quad400\,000;\quad200\,000\ \text{руб.}
\]
Соответственно проценты:
\[
40\,000;\quad30\,000;\quad20\,000;\quad10\,000\ \text{руб.}
\]
Добавляем одинаковые \(200\,000\) руб. долга:
\[
\begin{aligned}
A_1&=200\,000+40\,000=240\,000,\\
A_2&=200\,000+30\,000=230\,000,\\
A_3&=200\,000+20\,000=220\,000,\\
A_4&=200\,000+10\,000=210\,000.
\end{aligned}
\]
Теперь сравним \emph{соседние} платежи:
\[
A_1-A_2=A_2-A_3=A_3-A_4=10\,000.
\]
Можно сразу увидеть общую закономерность:
\[
A_j-A_{j+1}=d\cdot0{,}05
=200\,000\cdot0{,}05=10\,000.
\]

\textbf{Проверка.} От первого до четвёртого платежа три
одинаковых уменьшения:
\(3\cdot10\,000=30\,000\) руб., и действительно
\(240\,000-210\,000=30\,000\) руб.

\result\ Каждый следующий платёж меньше предыдущего на
\textbf{10\,000 руб.} Ваш числовой ответ верен; недоставало доказательства.

\textbf{Задача на закрепление.} Кредит 600\,000 руб. взят на четыре
года под 8\% годовых. Основной долг погашается одинаковыми частями.
Докажите, на сколько рублей каждый следующий ежегодный платёж
меньше предыдущего.

\clearpage
\tasktitle{9}
{\textbf{Условие.}} Кредит 1\,000\,000 руб. взят на пять лет с
равным ежегодным погашением основного долга. Первый платёж
не превышает 277\,000 руб., последний составляет не менее
214\,000 руб. Годовая ставка в процентах --- \emph{положительное
целое число}. Найдите ставку и общую сумму выплат.

\textbf{Ваше решение.} Верно составлены неравенства для первого
и последнего платежей, получены границы \(0{,}07\le p\le0{,}077\).
В записи рассматриваются «7\% и 7{,}7\%?». При вычислении итоговой
выплаты используется \(p=0{,}07\), и получено 1\,210\,000 руб.

\textbf{Ваш ответ.} В выборе ставки оставлен лишний вариант;
общая выплата 1\,210\,000 руб. верна.

\wrong\ Не учтено, что процентная ставка по условию должна быть
\emph{целым числом процентов}.

\textbf{Где именно возникает ошибка.} Последний правильный шаг ---
границы
\[
0{,}07\le p\le0{,}077.
\]
Первый незавершённый переход --- перечисление возможных ставок:
\(7{,}7\%\) попадает в числовой интервал, но не удовлетворяет
требованию целочисленности.

\textbf{Почему это неверно.} \(p\) --- доля единицы, поэтому
\(100p\) --- ставка в процентах. Условие требует, чтобы именно
\(100p\) было положительным целым числом. Из диапазона
от 7\% до 7{,}7\% подходит только 7\%.

\idea\ В задачах с целочисленной ставкой сначала найдите
допустимый интервал, затем отберите из него целые числа процентов.

\textbf{Исправление и полное решение.}
Ежегодно погашают одинаковую часть тела кредита:
\[
d=\frac{1\,000\,000}{5}=200\,000\ \text{руб.}
\]
Обозначим через \(p\) ставку как долю единицы.
Первый платёж начисляется на полный миллион:
\[
200\,000+1\,000\,000p\le277\,000
\quad\Longrightarrow\quad p\le0{,}077.
\]
Перед последним платежом остаётся 200\,000 руб.:
\[
200\,000+200\,000p\ge214\,000
\quad\Longrightarrow\quad p\ge0{,}07.
\]
Следовательно,
\[
0{,}07\le p\le0{,}077
\quad\Longleftrightarrow\quad
7\le100p\le7{,}7.
\]
Поскольку \(100p\) --- положительное целое число,
\[
\boxed{100p=7},\qquad \boxed{p=0{,}07}.
\]
Остатки долга в начале пяти лет (в тысячах рублей):
\(1000,800,600,400,200\).
Суммарно проценты начисляются на
\[
1000+800+600+400+200=3000
\quad\text{(тыс. руб.).}
\]
Общая выплата:
\[
1\,000\,000+0{,}07\cdot3\,000\,000
=\boxed{1\,210\,000\ \text{руб.}}
\]

\textbf{Проверка.} При ставке 7\% первый платёж
\(200\,000+70\,000=270\,000\le277\,000\) руб.;
последний
\(200\,000+14\,000=214\,000\ge214\,000\) руб.
Оба ограничения выполнены.

\result\ Ставка --- \textbf{7\% годовых};
общая выплата --- \textbf{1\,210\,000 руб.}

\textbf{Задача на закрепление.} Кредит 1\,200\,000 руб. выдан на
четыре года с равным ежегодным погашением основного долга.
Первый платёж не превышает 398\,000 руб., последний составляет
не менее 324\,000 руб. Годовая ставка в процентах ---
положительное целое число. Найдите ставку и общую сумму выплат.

\end{document}
